% ==============================================================================
% PHỤ LỤC C: CẤU HÌNH SIÊU THAM SỐ VÀ DỮ LIỆU THỰC NGHIỆM
% ==============================================================================

\chapter{CẤU HÌNH SIÊU THAM SỐ VÀ DỮ LIỆU THỰC NGHIỆM}
\label{app:hyperparams}

\section{Bảng tổng hợp cấu hình siêu tham số mô hình}

Chi tiết các thông số kỹ thuật được thiết lập trong quá trình tinh chỉnh mô hình thị giác thực phẩm và cấu hình tham số giải mã ngôn ngữ trên máy chủ Ollama được tổng hợp tại \cref{tab:hyperparams}.

\begin{table}[htbp]
\centering
\small
\caption{Bảng thiết lập các siêu tham số huấn luyện LoRA và suy luận mô hình}
\label{tab:hyperparams}
\begin{tabularx}{\textwidth}{>{\raggedright\arraybackslash}p{4.3cm} >{\raggedright\arraybackslash}p{3.6cm} >{\raggedright\arraybackslash}X}
    \toprule
    \textbf{Tham số cấu hình} & \textbf{Giá trị thiết lập} & \textbf{Ý nghĩa và mục đích kỹ thuật} \\
    \midrule
    Mô hình thị giác nền tảng & Qwen2-VL-7B-bnb-4bit & Mô hình đa phương thức lượng tử hóa 4-bit NormalFloat \\
    \addlinespace[3pt]
    Phương pháp tinh chỉnh & LoRA & Thích nghi ma trận hạng thấp tiết kiệm bộ nhớ đồ họa \\
    \addlinespace[3pt]
    Hạng phân rã ma trận $r$ & 16 & Cân bằng giữa độ phức tạp tham số và khả năng học đặc trưng \\
    \addlinespace[3pt]
    Hệ số khuếch đại $\alpha$ & 16 & Hệ số tỷ lệ chuẩn hóa tốc độ cập nhật của nhánh LoRA \\
    \addlinespace[3pt]
    Tỷ lệ rơi ngẫu nhiên & 0.0 & Tối ưu hóa tốc độ tính toán gradient với nhân Unsloth \\
    \addlinespace[3pt]
    Tốc độ học cơ sở & $2 \times 10^{-4}$ & Bước cập nhật trọng số trong không gian đạo hàm \\
    \addlinespace[3pt]
    Kích thước lô huấn luyện & 4 & Số lượng mẫu xử lý đồng thời trong mỗi bước cập nhật \\
    \addlinespace[3pt]
    Số chu kỳ huấn luyện & 5 & Số lượt duyệt qua toàn bộ tập dữ liệu ảnh thực phẩm \\
    \addlinespace[3pt]
    Nhiệt độ lấy mẫu trợ lý ảo & 0.25 & Giữ tính ổn định và chính xác cho các câu trả lời chuyên môn \\
    \addlinespace[3pt]
    Hệ số phạt lặp từ vựng & 1.18 & Ngăn chặn hiện tượng lặp lại cấu trúc câu khi sinh bài tập \\
    \addlinespace[3pt]
    Nhiệt độ lấy mẫu thị giác & 0.05 & Tối đa hóa tính tiền định khi xuất định dạng JSON \\
    \bottomrule
\end{tabularx}
\end{table}

\section{Mẫu dữ liệu phản hồi trích xuất dinh dưỡng định dạng JSON}

Dưới đây là mẫu phản hồi có cấu trúc trích xuất từ mô hình thị giác thực phẩm sau khi phân tích thành công ảnh đĩa cơm tấm sườn nướng trứng:

\begin{lstlisting}[language={}, caption={Mẫu đối tượng JSON trích xuất từ mô hình thị giác thực phẩm}, label={lst:json_sample}]
{
  "dish_name": "Com tam suon nuong trung op la",
  "portion_grams": 420,
  "confidence": 0.94,
  "cooking_method": "nuong va chien",
  "items": [
    {"name": "com tam", "grams": 200, "calories": 260, "carb": 56, "protein": 5, "fat": 1},
    {"name": "suon heo nuong", "grams": 130, "calories": 320, "carb": 2, "protein": 28, "fat": 22},
    {"name": "trung op la", "grams": 55, "calories": 90, "carb": 0.5, "protein": 6.5, "fat": 7},
    {"name": "dua leo ca chua", "grams": 35, "calories": 10, "carb": 2, "protein": 0.5, "fat": 0.1}
  ],
  "total_nutrition": {
    "calories": 680,
    "protein": 40.0,
    "carb": 60.5,
    "fat": 30.1
  }
}
\end{lstlisting}
